\section{Related Work}
\label{sec:rel_work}

Watermarking AI-generated content has become a critical area of research due to the rapid 
adoption of generative models and the concomitant concerns around provenance, copyright 
protection, and content authenticity. The proliferation of hyper-realistic synthetic media, 
including DeepFakes, has intensified the need for reliable attribution and verification 
mechanisms~\cite{gupta2024comprehensive, nguyen2022deep, verdoliva2020media}. 
Recent work spans techniques that embed imperceptible signatures into the image generation 
process itself, frequency-domain watermarking methods, and detection approaches that 
leverage deep learning for content authentication. These classes of methodologies are 
described next.

%% ----- SUBSECTION 1: Content Authentication and Detection Context -----
\noindent\textbf{Content Authentication in the Era of Generative AI.}
The growing sophistication of generative models has rendered traditional detection 
approaches increasingly insufficient. As surveyed by Gupta~et~al.~\cite{gupta2024comprehensive}, 
DeepFake detection methods have evolved from simple spatial artifact analysis to 
multi-modal fusion techniques combining spatial, temporal, and frequency-domain 
features. Frequency-based approaches have proven particularly effective: 
Neves~et~al.~\cite{neves2020ganprintr} demonstrated that the Discrete Fourier 
Transform can expose statistical artifacts in GAN-generated frames, compressing 
2D amplitude spectra into feature vectors for classification. 
Gu{\"e}ra and Delp~\cite{guera2018deepfake} showed that temporal inconsistencies 
introduced by face-swapping can be captured through CNN--RNN pipelines operating 
on frame-level features.
However, detection-based approaches are fundamentally reactive: they attempt to identify 
synthetic content after generation, without any cooperation from the generative model. 
This motivates \emph{proactive} approaches based on digital watermarking, which embed 
verifiable signatures during or after the generation process, enabling content attribution 
regardless of downstream analysis capabilities.

%% ----- SUBSECTION 2: Embedded Watermarking in Generative Models -----
\noindent\textbf{Embedded Watermarking in Generative Models.}
Rezaei~et~al.~\cite{rezaei2024lawa} propose LaWa, a watermarking framework that 
leverages the latent space of latent diffusion models (LDMs) to integrate watermarks 
during generation, achieving robustness across transformations while maintaining 
perceptual quality. This approach embeds signatures directly into generative pipelines 
rather than performing post-hoc modifications, demonstrating improved resilience and 
low false positive rates compared to earlier methods. A similar line of work focuses 
on embedding watermarks specifically within diffusion processes. 
Chen~et~al.~\cite{chen2025tagwm} introduce Dynamic Watermarks for diffusion models, 
using multi-stage embedding that includes a fixed watermark in the model's learned noise 
distribution and a dynamic, image-adapted watermark decoded via a fine-tuned neural 
decoder. This enables robust source verification with minimal visual impact.

Tree-Ring Watermarking~\cite{wen2023tree_ring} embeds concentric ring patterns into 
the Fourier space of the initial noise vector, producing watermarks that are invisible 
and structurally invariant to several geometric transformations. 
The Stable Signature~\cite{fernandez2023stable} fine-tunes the latent decoder of 
diffusion models to embed a recoverable binary signature into all generated outputs. 
More recently, methods such as ZoDiac~\cite{zhang2024attack}, 
Gaussian Shading~\cite{yang2024gaussian}, and ROBIN~\cite{huang2024robin} have 
introduced further innovations in latent-space watermark injection, adversarial concealment, and distribution-preserving embedding.

Despite these advances, existing in-generation methods exhibit important limitations. 
First, approaches relying on handcrafted detection metrics such as frequency-domain 
correlation or distance thresholds suffer from \emph{threshold instability}: the optimal 
decision boundary varies substantially across different image transformations. 
Second, most methods employ isotropic frequency structures (concentric rings, uniform 
grids, or random masks) that distribute watermark energy uniformly by radius, making 
them vulnerable to directional distortions including motion blur and anisotropic 
resampling. Third, existing verification strategies typically operate in a binary, 
threshold-dependent manner without producing calibrated confidence scores.

%% ----- SUBSECTION 3: Frequency-Domain Methods -----
\noindent\textbf{Frequency-Domain Watermarking: Fourier and Wavelet Transforms.}
The Discrete Fourier Transform (DFT) has long been a fundamental tool in signal and 
image processing. Several works~\cite{jia2023fourier, jiang2021focal, rao2021global, 
yu2022deep} have explored the use of DFT for image analysis, exploiting 
frequency-domain representations to enhance model performance and gain deeper 
insights into image characteristics. For instance, Baig~\cite{baig2022dft} employs 
DFT to globally assess the quality of blurred images, while 
Rao~\cite{rao2021global} uses frequency analysis to capture global image properties. 
These studies demonstrate that DFT is particularly effective for extracting global 
information from images, motivating its use in watermark injection and verification.

Concerning DWT-based approaches, the Discrete Wavelet Transform processes signals 
by decomposing them into multiple frequency bands at varying spatial resolutions. 
This multi-scale decomposition makes DWT especially well suited for capturing 
localized details. Prior studies~\cite{mohideen2008image, pimpalkhute2021digital} have 
successfully applied DWT to image denoising tasks. In particular, 
Tian~\cite{tian2023multi} highlights the benefit of DWT in jointly representing 
spatial and frequency information through multi-resolution analysis.

In the context of watermarking, classical approaches~\cite{barni2001improved,raval2003discrete,shensa2002discrete,tao2004robust} often exploit DWT subbands for embedding 
data. However, these methods were designed for natural images and do not account for 
the unique properties of diffusion-generated content, where the watermark can be 
injected directly into the generative noise.

%% ----- SUBSECTION 4: Deep Learning Watermarking -----
\noindent\textbf{Deep Learning for Watermarking and Steganography.}
Digital watermarking is a technique for safeguarding digital content by asserting 
copyright ownership. Its key distinction from steganography lies in the embedding 
priority: watermarking focuses on robustness, ensuring the embedded information 
remains detectable even after modifications or distortions, whereas steganography 
emphasizes invisibility. With the recent advances in deep learning, a plethora of 
learned watermarking methods have emerged. To the best of our knowledge, the 
HiDDeN~\cite{zhu2018hidden} approach was the first end-to-end deep watermarking 
framework that embeds resilient messages via a differentiable noise layer. 
More recently, Zhao~et~al.~\cite{zhao2024invisible} demonstrated that invisible 
image watermarks are provably removable using generative AI, underscoring the 
vulnerability of post-hoc watermarking methods and motivating in-generation 
approaches such as SpiderMark.

%% ----- SUBSECTION 5: Positioning of SpiderMark -----
\noindent\textbf{Positioning of SpiderMark.}
While many recent contributions embed watermarks within generative models, most focus 
on latent representation manipulation (e.g., LaWa), dynamic multi-stage embedding 
(e.g., Dynamic Watermarks), or distribution-preserving strategies (e.g., Gaussian 
Shading) for diffusion models, SpiderMark differentiates itself by proposing a truly 
hybrid watermarking framework that couples algorithmic frequency-domain injection with 
learning-based probabilistic verification. Unlike prior approaches that rely on 
handcrafted metrics or isotropic frequency patterns, SpiderMark introduces a structured 
spider-web pattern providing directional robustness and trains a lightweight verifier 
that outputs calibrated confidence scores, addressing the threshold instability that 
limits existing methods in real-world deployment scenarios.